\documentclass[11pt,letterpaper]{article}
\usepackage[T1]{fontenc}
\usepackage{lmodern,microtype,amsmath,amssymb,amsthm}
\usepackage[margin=1in]{geometry}
\usepackage[colorlinks=true,allcolors=blue!45!black]{hyperref}
\usepackage{xcolor}
\newtheorem{theorem}{Theorem}
\setlength{\emergencystretch}{2em}
\title{The missing frequencies of five-point prime variances}
\author{}
\date{November 27, 2026}
\begin{document}
\maketitle
\begin{abstract}
A four-point variance on consecutive observations misses at most one positive logarithmic frequency. We classify the five-point case. If the lag-three autocorrelation is nonzero, a quadratic equation and two phase tests give the complete set of missing frequencies. Its size is at most two, and shapes attaining two are dense in an explicit local family. If the lag-three autocorrelation vanishes, the count is either zero or infinite. Each missing frequency describes a particular oscillation of a positive counting measure invisible to the variance.
\end{abstract}

Write $P(z)=\sum_{j=0}^4h_jz^j$ for a nonzero real filter with $P(1)=P'(1)=0$ acting at five consecutive observation times. The inverse theorem in the ProveIt research note~\cite{prior} uses its symbol
\[
 D(\xi)=A+B2^{i\xi}+C3^{i\xi}+D_4\,4^{i\xi},\qquad
 (A,B,C,D_4)=(b_1,b_2,b_3,b_4),\quad b_k=k^3\sum_{j=0}^{4-k}h_jh_{j+k}.
\]
The notation $D_4$ distinguishes the last coefficient from the symbol. For the normalized spectral measure considered there, the variance response to a perturbation $\cos(\xi\log x)$ of the leading density is proportional to $D(-\xi)$ times a nonvanishing Gamma factor. Thus its real zeros give the invisible logarithmic frequencies. The zeta interpretation uses the explicit formula in the source article. All the following classifications of filters are unconditional.

\begin{theorem}[All shapes]
Set $\alpha=\log3/\log2$. If $C=0$, there are infinitely many positive invisible frequencies exactly when $A+Bz+D_4z^2$ has a root on the unit circle, and none otherwise. If $C\ne0$, solve
\[
 4AD_4 X^2+2B(A+D_4)X+(A-D_4)^2+B^2-C^2=0.
 \tag{1}\label{eq:poly}
\]
For each distinct real root $-1<X<1$, put $z=X+i\sqrt{1-X^2}=e^{i\theta}$ with $0<\theta<\pi$ and $w=-(A+Bz+D_4z^2)/C=e^{i\phi}$. Count that root once exactly when
\[
 (\phi-\alpha\theta)/(2\pi)\in\mathbb Z+\alpha\mathbb Z. \tag{2}\label{eq:group}
\]
The resulting count is the exact number of missing positive frequencies, at most two. If $\phi-\alpha\theta=2\pi(\alpha m-n)$, the corresponding frequency is $|\theta+2\pi m|/\log 2$.
\end{theorem}
\begin{proof}
The sum of the $h_j$ and their first moment vanish. Expanding $\sum_{ij}h_ih_j(i-j)^2$ gives $A+B/2+C/3+D_4/4=0$. Also $D(0)>0$ since $D(0)/6=\pi^{-1}\int_0^\infty|P(e^{ix})|^2 x^{-4}\,dx$ by the normalization in~\cite{prior}. When $C=0$ the response is periodic in $\xi$, proving the first assertion.

When $C\ne0$ let $z=2^{i\xi}$, $w=3^{i\xi}$. The equality $D(\xi)=0$ determines $w$ from $z$; the condition $|w|=1$ expands to (1), with $X=\Re z$. This quadratic cannot vanish identically. If $|A+Bz+D_4z^2|$ were constant $|C|$ on the unit circle, the polynomial would have only one nonzero coefficient. Indeed, the autocorrelation at the difference between its highest and lowest degrees is the product of the two extremal coefficients. If this difference is positive, that autocorrelation is nonzero, contrary to constant modulus. For a monomial the preceding linear relation forces the coefficient and $C$ to have distinct magnitudes, since the associated distance is $1,2$ or $4$, distinct from $3$.

If $z$ were real, $w$ would also be real, making $\xi\log2$ and $\xi\log3$ integer multiples of $\pi$. Irrationality of $\alpha$ forces $\xi=0$, excluded by $D(0)>0$. For each interior root the two choices of $z$ give conjugate phase configurations. Simultaneous lifts of the upper configuration require $\xi\log2=\theta+2\pi m$ and $\xi\log3=\phi+2\pi n$. These are compatible precisely in (2). Irrationality gives at most one pair of integers for each choice of $\phi$; conjugation accounts for the negative frequency. Distinct roots give distinct positive frequencies.
\end{proof}

\begin{theorem}[Attainment of the two-frequency bound]\label{thm:two}
Consider $P_{s,t}(z)=(1-z)^2(1+sz+tz^2)$. In an open neighborhood of $(s,t)=(8,14)$ every filter has five nonzero coefficients and nonzero lag-three autocorrelation. The shapes missing exactly two positive frequencies form a countable dense subset there. Shapes missing exactly one are also dense, and almost every shape misses none.
\end{theorem}
\begin{proof}
At the center the filter has coefficients $(1,6,-1,-20,14)$. Its symbol has coefficients $(-260,-1080,1728,896)$. The two roots in (1) lie respectively in $(-0.895,-0.894)$ and $(-0.580,-0.579)$. In a neighborhood of the center both roots and the phases $\theta_j,\phi_j$ of their upper configurations have real analytic local branches. Consider the map
\[
 (s,t)\longmapsto(\phi_1-\alpha\theta_1,\phi_2-\alpha\theta_2).
\]
Its determinant lies between $-0.008$ and $-0.007$ at the center. Here are formulas and bounds to verify this numerical inequality by rational arithmetic. For phase vectors write $[u,v]=\Re u\Im v-\Im u\Re v$. At each solution put
\[
 U=i(Bz+2D_4z^2),\quad V=iCw,\quad
 Q=A'+B'z+C'w+D_4'z^2,
\]
where a prime denotes differentiation in one parameter. Then $U\theta'+V\phi'=-Q$. Consequently the entry of the Jacobian is
\[
 \phi'-\alpha\theta'=\frac{-[U,Q]+\alpha[Q,V]}{[U,V]}.
\]
The coefficient derivatives in order $s,t$ are $(41,-352,405,0)$ and $(-60,16,108,64)$. Interval substitution in these formulas gives the following rational enclosures, with the roots listed in increasing order:
\[
 \begin{pmatrix}
 (0.0405,0.0408)&(-0.129,-0.128)\\
 (-0.0161,-0.0158)&(-0.144,-0.142)
 \end{pmatrix}.
\]
For reproducibility the script \texttt{check\_five\_nodes.py} encloses the square roots between consecutive multiples of $10^{-12}$ using integer square roots, bounds $1.5849<\alpha<1.5850$ by integer powers, and checks the determinant directly. The denominator $[U,V]$ is nonzero in both rows.
The inverse function theorem now identifies an open neighborhood of the center diffeomorphically with an open set of phase mismatches. The subgroup $G=2\pi(\mathbb Z+\alpha\mathbb Z)$ is countable and dense. Its product $G^2$ cuts out a countable dense set in every open rectangle; these parameters pass both tests. The points with exactly one coordinate in $G$ form a dense set, and the union of the two sets of coordinate lines has Lebesgue measure zero, also after a local smooth change of variables. All five coefficients remain nonzero after shrinking the neighborhood.
\end{proof}

In particular, a genuine five-node variance can converge on a spectrum with two independent adjustable amplitudes of invisible oscillation. If $\omega_1,\omega_2$ are its missing frequencies, one may take the positive leading density to be proportional to
$1+a\cos(\omega_1\log x)+b\cos(\omega_2\log x)$ with $|a|+|b|<1$. Substitution into the variance cancels both terms. The counting statistic retains both, with respective factors $(1+i\omega_j)^{-1}$ in the complex representation. The construction establishes distinct positive frequencies. The arithmetic nature of their ratio is a further question.

The dyadic example in the source article, recalled in \cite{prior}, has squared transfer function $|1-z|^4+|1-z^2|^4/8$ on the unit circle. By dilation, its symbol is the three-node second-difference symbol times $1+2^{i\xi}$, so it misses precisely the odd multiples of $\pi/\log 2$. The existence of a five-node example missing infinitely many frequencies already follows there. The two-frequency case and the criterion for all five-node shapes address the subsequent question in \cite{prior}. Greenhall studies a related spectral ambiguity for the Allan variance \cite{Greenhall}. A search of the repository and of literature on finite-difference variance ambiguity found no earlier statement of the two-frequency classification.

\begin{thebibliography}{1}
\bibitem{prior} ProveIt Contributors. \emph{Which four-point prime variances determine the multiplicities of zeta zeros?} November 27, 2026. See the finite family theorem, its positive countermodel, and the closing question on five consecutive observations. \href{https://github.com/VladimirReshetnikov/ProveIt/tree/00cadb961bbe77c74e5d082c4948f8884d58a970/Analysis/GeneralizedRiemannHypothesis/Research/grh-two-open-questions/zero-stencils}{Pinned source}.
\bibitem{Greenhall} C. A. Greenhall. Spectral ambiguity of Allan variance. \emph{IEEE Trans. Instrum. Meas.} \textbf{47} (1998), 623--627. \href{https://doi.org/10.1109/19.744312}{Published paper}.
\end{thebibliography}
\end{document}
